% Generated by experiments/scripts/export_latex.py -- do not edit by hand.
% Source: experiments/data/2-processed-tables-and-evaluations/revised-usd-12m/same-window/eval_summary.json
% Source: experiments/data/2-processed-tables-and-evaluations/revised-usd-12m/holdout-w0/eval_summary.json
% Source: experiments/data/2-processed-tables-and-evaluations/revised-usd-12m/window-w1/eval_summary.json
\begin{table}[htbp]
\centering
\caption[AWP 옆에 놓은, 고르게 다시 출발하는 엔도스랭크와 활동 재시작을 쓰는 엔도스랭크.]{엔도스랭크를 고르게 다시 출발하게 한 것과, 허락하는 활동에 따라 무게를 단 활동 재시작을 쓰게 한 것을 AWP 옆에 놓았어요. 활동 재시작은 AWP의 규칙을 허락 그래프에 쓴 것이에요. 매칭 코호트(계약 $n=6{,}483$개)의 지표 가족 평균 켄달의 $\tau$를 보여 줘요. 또 두 스펜더 라벨과의 $\tau$를 W0에서(스펜더 계약 $n=7{,}914$개, 점수는 2026년 3월 31일에 동결), 그리고 W1에서($n=7{,}960$, 2026년 6월 30일에 동결) 보여 줘요. 두 엔도스랭크 행 모두 각 마지막 허락에 $\sigma(\Delta t)V(z)$만큼 무게를 줘요. 신뢰구간(CI)은 계약을 다시 뽑아 구한 95\% 백분위 부트스트랩 구간이에요(짝지은 다시 뽑기 400번, 시드 42).}
\label{tab:endorserank-restarts}
\fitwidth{%
\begin{tabular}{llccccccc}
\toprule
 & & \multicolumn{3}{c}{같은 기간, 매칭 코호트} & \multicolumn{2}{c}{W0, 스펜더} & \multicolumn{2}{c}{W1, 스펜더} \\
\cmidrule(lr){3-5} \cmidrule(lr){6-7} \cmidrule(lr){8-9}
방법 & 다시 출발 & 송금 & 허락 & 시빌 & 새 허락 쌍 & 새 송금자 & 새 허락 쌍 & 새 송금자 \\
\midrule
EndorseRank & 고르게 & \begin{tabular}[c]{@{}c@{}}0.196\\[-1pt]{\footnotesize [0.180, 0.216]}\end{tabular} & \begin{tabular}[c]{@{}c@{}}0.560\\[-1pt]{\footnotesize [0.548, 0.571]}\end{tabular} & \begin{tabular}[c]{@{}c@{}}0.026\\[-1pt]{\footnotesize [0.012, 0.041]}\end{tabular} & \begin{tabular}[c]{@{}c@{}}0.303\\[-1pt]{\footnotesize [0.287, 0.319]}\end{tabular} & \begin{tabular}[c]{@{}c@{}}0.272\\[-1pt]{\footnotesize [0.257, 0.286]}\end{tabular} & \begin{tabular}[c]{@{}c@{}}0.299\\[-1pt]{\footnotesize [0.284, 0.314]}\end{tabular} & \begin{tabular}[c]{@{}c@{}}0.259\\[-1pt]{\footnotesize [0.243, 0.275]}\end{tabular} \\
EndorseRank & 허락하는 활동에 따라 & \begin{tabular}[c]{@{}c@{}}0.251\\[-1pt]{\footnotesize [0.233, 0.271]}\end{tabular} & \begin{tabular}[c]{@{}c@{}}0.601\\[-1pt]{\footnotesize [0.590, 0.613]}\end{tabular} & \begin{tabular}[c]{@{}c@{}}0.056\\[-1pt]{\footnotesize [0.042, 0.072]}\end{tabular} & \begin{tabular}[c]{@{}c@{}}0.332\\[-1pt]{\footnotesize [0.315, 0.347]}\end{tabular} & \begin{tabular}[c]{@{}c@{}}0.270\\[-1pt]{\footnotesize [0.256, 0.285]}\end{tabular} & \begin{tabular}[c]{@{}c@{}}0.335\\[-1pt]{\footnotesize [0.319, 0.349]}\end{tabular} & \begin{tabular}[c]{@{}c@{}}0.236\\[-1pt]{\footnotesize [0.219, 0.251]}\end{tabular} \\
AWP & 송금하는 활동에 따라 & \begin{tabular}[c]{@{}c@{}}0.736\\[-1pt]{\footnotesize [0.728, 0.744]}\end{tabular} & \begin{tabular}[c]{@{}c@{}}0.226\\[-1pt]{\footnotesize [0.210, 0.245]}\end{tabular} & \begin{tabular}[c]{@{}c@{}}0.419\\[-1pt]{\footnotesize [0.408, 0.434]}\end{tabular} & \begin{tabular}[c]{@{}c@{}}0.196\\[-1pt]{\footnotesize [0.178, 0.213]}\end{tabular} & \begin{tabular}[c]{@{}c@{}}0.409\\[-1pt]{\footnotesize [0.396, 0.421]}\end{tabular} & \begin{tabular}[c]{@{}c@{}}0.157\\[-1pt]{\footnotesize [0.137, 0.178]}\end{tabular} & \begin{tabular}[c]{@{}c@{}}0.398\\[-1pt]{\footnotesize [0.386, 0.412]}\end{tabular} \\
\bottomrule
\end{tabular}
}
\end{table}
